% Generated draft: requires mathematical and source review.
\documentclass[11pt]{article}
\usepackage[margin=25mm]{geometry}
\usepackage{fontspec}
\setmainfont{Latin Modern Roman}
\usepackage{amsmath,amssymb,mathtools,unicode-math}
\setmathfont{Latin Modern Math}
\usepackage{enumitem,needspace,xurl,hyperref}
\setlength{\parindent}{0pt}
\setlength{\parskip}{5pt}
\setlength{\emergencystretch}{3em}
\newcommand{\N}{\mathbb{N}}
\newcommand{\Z}{\mathbb{Z}}
\newcommand{\Q}{\mathbb{Q}}
\newcommand{\R}{\mathbb{R}}
\newcommand{\C}{\mathbb{C}}
\newcommand{\F}{\mathbb{F}}
\newcommand{\A}{\mathbb{A}}
\newcommand{\PP}{\mathbb{P}}
\newcommand{\E}{\mathbb{E}}
\newcommand{\eps}{\varepsilon}
\newcommand{\dd}{\,\mathrm{d}}
\begin{document}
\section*{UnsolvedMath ID: 30002671}
\textbf{Source problem code:} catalogue entry
\subsection{Definitions and mathematical statement}
The supplied record describes a one-round Voronoi game in which the first player places \(t\) facilities and the second player places one facility. Here \(t\) denotes the first player's facility count; the admissible range of \(t\) is not specified. The record uses the notation \(VR_{t:1}\), but does not define the quantity it denotes.

The stated question is whether the bound
\[
VR_{t:1}\leq \frac{t-1}{t+1}
\]
is sharp for this game.

This formulation remains incomplete: the underlying space or class of instances, distance rule, legal facility placements, allocation of points equidistant from facilities, payoff, and optimization defining \(VR_{t:1}\) are not supplied. Consequently, neither the precise quantified inequality nor the meaning of sharpness can be fixed from this record. In particular, it is unspecified whether sharpness means equality for an admissible instance or approximation arbitrarily closely, and whether it is sought for every admissible \(t\). These qualifications require resolution before a complete mathematical statement can be given.
\subsection{Short English statement}
The question asks whether the stated upper bound \((t-1)/(t+1)\) for \(VR_{t:1}\) in a one-round Voronoi game can be improved. The game rules, the definition of \(VR_{t:1}\), the range of \(t\), and the intended meaning of sharpness are missing.
\subsection{Sources}
\begin{itemize}
\item Supplied local UnsolvedMath registry record, permanent ID \(30002671\), source problem code \texttt{catalogue entry}, titled ``Sharpness of the One-Round Voronoi Game Bound,'' in the Oberwolfach Reports Open Problems collection. No individual proposer is identified; the supplied proposed year is 2014.
\item Supplied source citation: \emph{Discrete Geometry (2015)}, Oberwolfach Reports, DOI 10.4171/owr/2014/40, pp.\ 2235--2296. \url{https://doi.org/10.4171/owr/2014/40}. The background describes these reports as records of mathematical workshops and their open-problem discussions. The relationship between the citation's 2015 label and the supplied 2014 year is unresolved.
\item Supplied literature reference: Gerbner et al., ``advantage in the discrete Voronoi game.'' \url{https://doi.org/10.7155/jgaa.00331}.
\item Supplied literature reference: ``Oberwolfach report.'' \url{https://ems.press/content/serial-article-files/46530}.
\item Literature assessment dated exactly \texttt{2026-08-21}: the supplied curation labels the problem open and reports that discrete Voronoi-game bounds are established while exact sharpness of the stated bound remains open in the cited primary literature and the curation's search. This is supplied dated curation derived from the Oberwolfach material, not an independently verified assessment. The record explicitly calls for independent verification of time-sensitive status and cited claims.
\end{itemize}
\end{document}
